\chapter{ROS 2 y PX4 para principiantes}

\begin{enpocas}
La terminal en cinco minutos, qué es ROS~2 (nodos, tópicos, mensajes, parámetros, launch), qué es PX4,
y cómo se hablan los dos por medio del agente uXRCE-DDS.
\end{enpocas}

\section{La terminal en cinco minutos}

\begin{lstlisting}[style=bash]
pwd                      # ¿en qué carpeta estoy?
ls                       # ¿qué hay aquí?
cd ~/ros2_ws             # entrar a una carpeta
cat archivo.txt          # ver un archivo
Ctrl+C                   # detener el programa que está corriendo
\end{lstlisting}

\begin{consejo}
Tab completa nombres y la flecha arriba recupera la orden anterior.
\end{consejo}

\section{ROS 2}

\textbf{ROS~2} (\emph{Robot Operating System}) conecta muchos programas pequeños. Cada programa es un
\textbf{nodo}; los nodos se mandan \textbf{mensajes} por canales con nombre llamados \textbf{tópicos}.
Uno \emph{publica}, otros se \emph{suscriben}, sin conocerse entre sí.

\begin{ejemplo}
Una cocina de restaurante: quien toma pedidos los pega en un tablero (tópico), el cocinero los toma
de ahí y deja los platos en la barra (otro tópico), de donde los recoge el mesero. Nadie habla con
nadie directamente.
\end{ejemplo}

En este proyecto:
\begin{longtable}{>{\raggedright\arraybackslash}p{4.5cm}p{10.5cm}}
\toprule
\textbf{Tópico} & \textbf{Qué lleva} \\ \midrule \endhead
\cod{/camera} & imágenes de la cámara del dron (Gazebo $\rightarrow$ ROS) \\
\cod{/detecciones} & recuadros de personas que encuentra YOLO \\
\cod{/personas/estimaciones} & una posición en el suelo por cada detección \\
\cod{/personas/candidatas} & candidatas fusionadas con su confianza y número de detecciones \\
\cod{/personas/gps}, \cod{/personas/alerta} & aviso de cada persona confirmada \\
\cod{/fmu/out/vehicle_odometry} & posición y orientación estimadas por PX4 \\
\cod{/fmu/in/trajectory_setpoint} & consigna de posición que manda la misión a PX4 \\
\cod{/x500/lidar/points} & nube de puntos del LiDAR \\
\bottomrule
\end{longtable}

\begin{lstlisting}[style=bash]
ros2 node list
ros2 topic list
ros2 topic echo /personas/alerta
ros2 topic hz /camera
\end{lstlisting}

Los ajustes de cada nodo son \textbf{parámetros}; un archivo \textbf{launch} arranca todos los nodos con
sus parámetros. El nuestro, \cod{simulacion.launch.py}, también genera el bosque y arranca Gazebo,
PX4 o ArduPilot y MAVProxy.

\section{PX4}

\textbf{PX4} es uno de los dos pilotos automáticos de código abierto más usados (el otro es ArduPilot).
Estima el estado con un filtro de Kalman (EKF2), controla en cascada (posición $\rightarrow$
velocidad $\rightarrow$ actitud $\rightarrow$ velocidad angular) y tiene modos de vuelo. Dos son clave
aquí:
\begin{description}
  \item[Offboard:] PX4 obedece consignas de una computadora externa (nuestra misión) que debe llegar
        al menos dos veces por segundo; si dejan de llegar, PX4 se protege.
  \item[Neural Control] (módulo \cod{mc_nn_control}, en la versión 1.17): una red neuronal calcula
        directamente los comandos de los cuatro motores (capítulo~\ref{cap:control-neuronal}).
\end{description}

\section{El puente: uXRCE-DDS}

PX4 tiene su propio sistema interno de mensajes (uORB). El \textbf{agente Micro XRCE-DDS} los traduce a
tópicos de ROS~2 con el prefijo \cod{/fmu/}: \cod{/fmu/out/...} lo que publica PX4,
\cod{/fmu/in/...} lo que recibe. Los tipos de mensaje vienen del paquete \cod{px4_msgs}, que debe ser
de la misma versión que PX4 (\cod{release/1.17}).

\begin{leccion}
En PX4~1.17 algunos mensajes son \emph{versionados}: el estado del vehículo se publica en
\cod{/fmu/out/vehicle_status_v1}, no en \cod{/fmu/out/vehicle_status} como decía el diseño original.
Escuchar el tópico equivocado no da error: simplemente nunca llega nada.
\end{leccion}

\section{Workspace y compilación}

\begin{lstlisting}[style=bash]
cd ~/ros2_ws
colcon build --symlink-install
source install/setup.bash
\end{lstlisting}
El workspace contiene \cod{mision} (Python), \cod{mision_gz_plugins} (C++), \cod{px4_msgs} y, para
ArduPilot, \cod{ardupilot_gazebo}.
